\subsubsection{Symbian OS}

\sourcenotebox{\textbf{Source Note:} The listings below are taken from the
local Symbian EPOC Kernel Architecture 2 (EKA2) source tree in this workspace
under \texttt{kernel/eka/}.
The upstream repository is
\url{https://github.com/SymbianSource/oss.FCL.sf.os.kernelhwsrv}. Each caption
names the relevant source file or files.}

Symbian EKA2 is best understood as a layered nanokernel-plus-kernel design.
The nanokernel exports the minimal execution substrate, while the larger kernel
executive builds processes, drivers, inter-process communication (IPC), memory
objects, and security policy on top of that substrate
\citep{sales2005, symbian_eka2_kernel}. This is not a pure microkernel because
drivers and file systems still run in kernel space, but it is also not a flat
monolith because the nanokernel boundary is explicit in the headers, boot
flow, and object model.

\paragraph{Architecture and Design}

The architectural split is visible in the nanokernel header itself.
Listing~\ref{lst:kernel-nkern-boundary} shows that the standalone nanokernel
build path is explicit and that the header surface is centered on low-level
facilities such as queues, deferred function calls (DFCs), tracing, atomics,
fault entry, and idle handling. This is the smallest execution substrate on
which the larger kernel is intended to depend.

\begin{listing}[H]
\begin{minted}{cpp}
#ifdef	__STANDALONE_NANOKERNEL__
#undef	__IN_KERNEL__
#define	__IN_KERNEL__
#endif

#include <e32const.h>
#include <nklib.h>
#include <dfcs.h>
#include <nk_trace.h>
#include <e32atomics.h>

extern "C" {
IMPORT_C void NKFault(const char* file, TInt line);
void NKIdle(TUint32 aStage);
}
\end{minted}
\caption{Standalone nanokernel boundary and low-level surface. Source:
\texttt{include/nkern/nkern.h}.}
\label{lst:kernel-nkern-boundary}
\end{listing}

The layering also shows up in initialization. Once the first thread's
nanokernel context is prepared, boot explicitly hands control to the memory
model and then the variant or application-specific integrated circuit (ASIC)
layer instead of fusing those concerns into a single monolithic startup
routine. Listing~\ref{lst:kernel-init-split} captures that handoff.

\begin{listing}[H]
\begin{minted}{cpp}
	info.iSlowExecTable = (const SSlowExecTable*)EpocSlowExecTable;
	info.iParameterBlock = (const TUint32*)aInfo;
	info.iParameterBlockSize = 0;
	info.iCpuAffinity = 0;
	info.iGroup = 0;

	TAny* ec = t->iNThread.iExtraContext;
	TInt ecs = t->iNThread.iExtraContextSize;
	NKern::Init(&t->iNThread, info);
	t->iNThread.SetExtraContext(ec, ecs);

	__KTRACE_OPT(KBOOT,DEBUGPRINT("NKern::CurrentThread() = %08x",
		NKern::CurrentThread()));

	M::Init2AP();
	A::Init2AP();
\end{minted}
\caption{Boot-time handoff from nanokernel setup into the memory model and
variant layer. Source: \texttt{kernel/sinit.cpp}.}
\label{lst:kernel-init-split}
\end{listing}

Compared with Linux, this is a more deliberately stratified kernel interior.
Compared with QNX, it still keeps too much functionality in kernel space to
provide strong component restartability.

\paragraph{Power and Mobile-First Behavior}

Symbian was designed for battery-constrained phones, and the idle path proves
it. Listing~\ref{lst:kernel-power} shows that when the scheduler runs out of
work, the kernel first asks a central processing unit (CPU) idle handler, then
the power model, and only then falls back to the ASIC-specific idle hook. The
power subsystem is therefore not an afterthought layered on top of scheduling;
scheduler quiescence and power management are treated as closely related
events.

\begin{listing}[H]
\begin{minted}{cpp}
extern "C" void NKIdle(TUint32 aStage)
	{
	SCpuIdleHandler* cih = NKern::CpuIdleHandler();
#ifdef __SMP__
	TSubScheduler& ss = SubScheduler();
	if (cih && cih->iHandler)
		(*cih->iHandler)(cih->iPtr, aStage, ss.iUncached);
#else
	if (cih && cih->iHandler)
		(*cih->iHandler)(cih->iPtr, aStage);
#endif
	else if (K::PowerModel)
		K::PowerModel->CpuIdle();
	else
		Arm::TheAsic->Idle();
	}

class DPowerManager : public DPowerModel
	{
public:
	void CpuIdle();
	void RegisterUserActivity(const TRawEvent& anEvent);
	TInt PowerHalFunction(TInt aFunction, TAny* a1, TAny* a2);
	void AbsoluteTimerExpired();
	void SystemTimeChanged(TInt anOldTime, TInt aNewTime);
	TSupplyStatus MachinePowerStatus();
	};
\end{minted}
\caption{Idle handling and power-model integration. Source:
\texttt{kernel/arm/cinit.cpp} and \texttt{kernel/power.cpp}.}
\label{lst:kernel-power}
\end{listing}

\paragraph{Drivers, Objects, and Extensibility}

Symbian keeps drivers in kernel space, but it still exposes a clear and fairly
systematic registration path. Listing~\ref{lst:kernel-drivers-objects} shows
three parts of that structure: logical devices are installed as kernel
objects, the kernel keeps shared global service queues for DFC and system work,
and core objects receive stable IDs at construction time. That is what makes
EKA2 a structured embedded monolith rather than a loose collection of
subsystems.

\begin{listing}[H]
\begin{minted}{cpp}
EXPORT_C TInt Kern::InstallLogicalDevice(DLogicalDevice* aDevice)
	{
	aDevice->SetProtection(DObject::EGlobal);
	aDevice->iObjectFlags |= DObject::EObjectExtraReference;
	TInt r=aDevice->Install();
	if (r==KErrNone)
		r=K::AddObject(aDevice,ELogicalDevice);
	return r;
	}

TMachineConfig* K::MachineConfig;
DThread* K::TheKernelThread;
DProcess* K::TheKernelProcess;
TDfcQue* K::DfcQ0;
TDfcQue* K::DfcQ1;
TDfcQue* K::TimerDfcQ;
DObjectCon* K::Containers[ENumObjectTypes];

EXPORT_C DObject::DObject()
	{
	iAccessCount=1;
	iObjectId = __e32_atomic_add_ord64(&NextObjectId, TUint64(1<<18));
	}
\end{minted}
\caption{Kernel driver installation, shared global state, and base object
construction. Source: \texttt{kernel/device.cpp}, \texttt{kernel/sglobals.cpp},
and \texttt{kernel/object.cpp}.}
\label{lst:kernel-drivers-objects}
\end{listing}

The benefit is low overhead and straightforward entry points. The cost is weak
runtime fault isolation: a driver bug is still a kernel bug.

\paragraph{User Space, System Calls, and Safe User Memory}

Symbian enforces the user-kernel boundary with per-process address spaces and a
shared super page. Listing~\ref{lst:kernel-userspace-boundary} shows both
sides: the super page makes selected global state intentionally visible in each
process, while executive-call tables and per-thread dispatch pointers keep the
system-call path explicit and inspectable instead of implicit or opaque.

\begin{listing}[H]
\begin{minted}{cpp}
struct SSuperPageBase
	{
	TInt		iRamDriveSize;
	TInt		iTotalRomSize;
	TInt		iTotalRamSize;
	TLinAddr	iPageDir;
	TInt		iKernelFault;
	TLinAddr	iRamBootData;
	TLinAddr	iRomBootData;
	TLinAddr	iBootTable;
	TLinAddr	iCodeBase;
	TUint32		iBootFlags;
	TLinAddr	iBootAlloc;
	TLinAddr	iMachineData;
	TUint		iActiveVariant;
	TLinAddr	iPrimaryEntry;
	TInt		iInitialHeapSize;
	TInt		iDebugPort;
	TUint		iCpuId;
	TLinAddr	iUncachedAddress;
	TLinAddr	iRootDirList;
	TInt		iHwStartupReason;
	};

#define	FAST_EXEC_BEGIN	\
	GLDEF_D const TUint32 EpocFastExecTable[] =		\
		{

#define	SLOW_EXEC_BEGIN	\
	GLDEF_D const TUint32 EpocSlowExecTable[] =		\
		{

	const SNThreadHandlers* iHandlers;
	const SFastExecTable* iFastExecTable;
	const SSlowExecEntry* iSlowExecTable;
	TLinAddr iSavedSP;
\end{minted}
\caption{Shared super-page state and explicit executive-call dispatch
machinery. Source: \texttt{include/kernel/kernboot.h},
\texttt{include/kernel/execs.h}, and \texttt{include/nkern/nk\_priv.h}.}
\label{lst:kernel-userspace-boundary}
\end{listing}

Because user pointers are unsafe by default, sensitive kernel-to-user crossings
are wrapped in trap harnesses and safe-copy helpers. Listing~\ref{lst:kernel-safe-user}
shows the ordinary and paging-aware exception traps together with a
\texttt{USafeWrite()}-based IPC callback path. Symbian still runs many services
in kernel mode, but it does not assume user memory can be touched recklessly.

\begin{listing}[H]
\begin{minted}{cpp}
class TExcTrap
	{
public:
	IMPORT_C TInt Trap();
	IMPORT_C TInt Trap(TExcTrapHandler aHandler);
	inline void UnTrap();
	IMPORT_C void Exception(TInt aResult);
public:
	TUint32 iState[EXC_TRAP_CTX_SZ];
	TExcTrapHandler iHandler;
	DThread* iThread;
	};

class TPagingExcTrap
	{
public:
	TInt Trap();
	inline void UnTrap();
	void Exception(TInt aResult);
public:
	TUint32 iState[EXC_TRAP_CTX_SZ];
	DThread* iThread;
	};

#define XTRAP_PAGING(_r,_s)	{								\
							TPagingExcTrap __t;				\
							if (((_r)=__t.Trap())==0)		\
								{							\
								_s;							\
								__t.UnTrap();				\
								}							\
							}

void TServerMessage::CallbackFunc(TAny* aData, TUserModeCallbackReason aReason)
	{
	TServerMessage* req = (TServerMessage*)aData;
	if (aReason == EUserModeCallbackRun && req->iResult == KErrNone)
		{
		RMessageU2 userData(*req->iMessageData);
		K::USafeWrite(req->iMessagePtr, &userData, sizeof(userData));
		}
	TClientRequest::CallbackFunc(aData, aReason);
	}
\end{minted}
\caption{Exception trapping and safe user-memory crossing. Source:
\texttt{include/kernel/kern\_priv.h} and \texttt{kernel/sipc.cpp}.}
\label{lst:kernel-safe-user}
\end{listing}

\paragraph{Interrupts, Deferred Work, and Timers}

Symbian uses a staged interrupt model. Listing~\ref{lst:kernel-interrupts}
shows three distinct pieces of that design: global interrupt-dispatch state in
the nanokernel, explicit DFC queue types and priorities, and timer completion
that can be requested in DFC context. The design is a sensible embedded
compromise: short work stays near interrupt context, while longer work is
pushed into DFC threads.

\begin{listing}[H]
\begin{minted}{cpp}
extern "C" {
SArmInterruptInfo ArmInterruptInfo;
}

TScheduler TheScheduler;

const TInt KNumDfcPriorities=8;
const TInt KMaxDfcPriority=KNumDfcPriorities-1;

class TDfcQue : public TPriList<TDfc,KNumDfcPriorities>
	{
public:
	IMPORT_C TDfcQue();
	inline TBool IsEmpty();
	static void ThreadFunction(TAny* aDfcQ);
public:
	NThreadBase* iThread;
	};

EXPORT_C TInt NTimer::OneShot(TInt aTime, TBool aDfc)
	{
	__NK_ASSERT_DEBUG(aTime>=0);
	__NK_ASSERT_DEBUG(iFunction!=NULL);
	TInt irq=NKern::DisableAllInterrupts();
	if (iState!=EIdle)
		{
		NKern::RestoreInterrupts(irq);
		return KErrInUse;
		}
	iCompleteInDfc=TUint8(aDfc?1:0);
	iTriggerTime=TheTimerQ.iMsCount+(TUint32)aTime;
	TheTimerQ.Add(this);
	NKern::RestoreInterrupts(irq);
	return KErrNone;
	}
\end{minted}
\caption{Interrupt state, deferred-call queues, and timer completion policy.
Source: \texttt{nkern/arm/ncglob.cpp}, \texttt{include/nkern/dfcs.h}, and
\texttt{nkern/nk\_timer.cpp}.}
\label{lst:kernel-interrupts}
\end{listing}

\paragraph{Platform Security and Capability Checks}

Symbian's platform security model is capability-based. The capability list is
defined in a shared header, and checks happen at the process or thread
boundary. Listing~\ref{lst:kernel-capabilities} shows both the shared
enumeration and the concrete process or thread check path in
\texttt{ssecure.cpp}. This is simpler than Linux's layered discretionary,
capability, and Linux Security Module (LSM)-based permission models, but it
fits a signed, vendor-controlled phone ecosystem well.

\begin{listing}[H]
\begin{minted}{cpp}
enum TCapability
	{
	ECapabilityTCB				= 0,
	ECapabilityCommDD			= 1,
	ECapabilityPowerMgmt		= 2,
	ECapabilityMultimediaDD		= 3,
	};

EXPORT_C TBool DProcess::DoHasCapability(TCapability aCapability,
	const char* aContextText)
	{
	if (HasCapabilityNoDiagnostic(aCapability))
		return ETrue;
	return PlatSec::CapabilityCheckFail(this,aCapability,aContextText)==KErrNone;
	}

EXPORT_C TBool DThread::DoHasCapability(TCapability aCapability,
	const char* aContextText)
	{
	if(iOwningProcess->HasCapabilityNoDiagnostic(aCapability))
		return ETrue;
	return PlatSec::CapabilityCheckFail(this,aCapability,aContextText)==KErrNone;
	}
\end{minted}
\caption{Capability definitions and kernel-side capability enforcement. Source:
\texttt{include/e32capability.h} and \texttt{kernel/ssecure.cpp}.}
\label{lst:kernel-capabilities}
\end{listing}

\paragraph{Fault Handling and Stability}

Fault isolation is where the limits of the design are most obvious. User-space
failures are isolated by the memory management unit (MMU) and safe-copy
scaffolding, but kernel-space failures remain catastrophic because drivers and
core services share the scheduler's address space. Listing~\ref{lst:kernel-faults}
shows how blunt the failure path is: the nanokernel exposes a direct fatal
fault entry point, and the thread state machine treats impossible internal
state as a kernel-fatal condition.

\begin{listing}[H]
\begin{minted}{cpp}
IMPORT_C void NKFault(const char* file, TInt line);
#define FAULT()		NKFault(__FILE__,__LINE__)

	__KTRACE_OPT(KTHREAD,Kern::Printf(
		"UNKNOWN THREAD MSTATE!! thread %O Mstate=%d, operation=%d, parameter=%d",
		aThread,aThread->iMState,anOperation,aParameter));
	Kern::Fault("THREAD STATE",aThread->iMState);
\end{minted}
\caption{Kernel-fatal fault paths for nanokernel and thread-state corruption.
Source: \texttt{include/nkern/nkern.h} and \texttt{kernel/sthread.cpp}.}
\label{lst:kernel-faults}
\end{listing}

That places Symbian between Linux and QNX. It has better internal layering
than a flat monolith, but it does not deliver the restartability or hard fault
containment of a true user-space driver architecture. For its era and target
hardware, that was a defensible trade-off: low overhead, explicit control, and
a kernel structure a phone manufacturer could reason about, test, and ship on
severely constrained devices \citep{sales2005, symbian_eka2_kernel}.
